\section{Caracterización de los cuadrípolos individuales}

Para cada cuadrípolo se implementaron cuatro configuraciones de medición. En ellas se excitó uno de los puertos y se midieron las tensiones de entrada, de salida y sobre las resistencias auxiliares. Estas mediciones permitieron obtener las corrientes necesarias para calcular los distintos parámetros.

\subsection{Esquemáticos de medición}

Cabe aclarar que, para todas las mediciones realizadas, se tuvo en cuenta la resistencia de la fuente, $R_f$, la cual se supuso conectada en serie y con un valor de $100\Omega$. Por lo tanto, la corriente que circula por el puerto conectado a la fuente se calculó siempre como $I=\frac{V_R}{R_f}$, donde $V_R$ es la caída de potencial sobre $R_f$.

Inicialmente se dejó el puerto 2 en circuito abierto, como se puede ver en la Figura \ref{fig:conf_1}. Con esta configuración se pudieron obtener los valores de $Z_{11}$, $Z_{21}$, $A$ y $C$.

\begin{figure}[H]
\centering

\begin{tikzpicture}[transform shape]
	% Paths, nodes and wires:
	\node[fourport, xscale=1.6, yscale=1.6] at (6.456, 2.353){};
	\node[fourport, xscale=0.3, yscale=0.3] at (5.523, 3.239){};
	\draw (2.25, 3.375) to[american resistor, l={$R_f$}] (5, 3.375);
	\draw (2.25, 3.375) to[vco] (2.25, 1.375);
	\draw (2.25, 1.375) -- (5, 1.375);
	\draw (8, 3.375) -- (9.75, 3.375);
	\draw (8, 1.375) -- (9.75, 1.375);
	\draw[-latex] (8.5, 1.625) -- (8.5, 3.125);
	\draw[-latex] (4.75, 1.625) -- (4.75, 3.125);
	\draw[-latex] (3, 2.875) -- (4, 2.875);
	\draw[-latex] (9.25, 3.625) -- (8.25, 3.625);
	\node[shape=rectangle, minimum width=0.465cm, minimum height=0.215cm](N1) at (8.75, 3.875){} node[anchor=center] at (N1.text){$I_2$};
	\node[shape=rectangle, minimum width=0.715cm, minimum height=0.465cm](N2) at (8.875, 2.5){} node[anchor=center] at (N2.text){$V_2$};
	\node[shape=rectangle, minimum width=0.465cm, minimum height=0.465cm](N3) at (3.5, 2.5){} node[anchor=center] at (N3.text){$I_1$};
	\node[shape=rectangle, minimum width=0.465cm, minimum height=0.465cm](N4) at (4.5, 2.375){} node[anchor=center] at (N4.text){$V_1$};
	\node[shape=rectangle, minimum width=0.715cm, minimum height=0.465cm](N5) at (1.375, 2.5){} node[anchor=center] at (N5.text){$V_S$};
\end{tikzpicture}

\caption{Primera configuración utilizada, generando un circuito abierto en el puerto 2.}
\label{fig:conf_1}
\end{figure}

Posteriormente se dejó el puerto 1 en circuito abierto, como se puede ver en la Figura \ref{fig:conf_2}. Con esta configuración se obtuvieron los valores de $Z_{22}$ y $Z_{12}$.

\begin{figure}[H]
\centering
\begin{tikzpicture}[transform shape]
	% Paths, nodes and wires:
	\node[fourport, xscale=1.6, yscale=1.6] at (6.544, 2.272){};
	\node[fourport, xscale=0.3, yscale=0.3] at (5.611, 3.158){};
	\draw (10.75, 3.25) to[american resistor, mirror, l={$R_f$}] (8, 3.25);
	\draw (10.75, 3.25) to[vco, mirror] (10.75, 1.25);
	\draw (8, 1.25) -- (10.75, 1.25);
	\draw (3.338, 3.169) -- (5.088, 3.169);
	\draw (3.338, 1.169) -- (5.088, 1.169);
	\draw[-latex] (8.25, 1.5) -- (8.25, 3);
	\draw[-latex] (4.838, 1.419) -- (4.838, 2.919);
	\draw[-latex] (3.588, 3.419) -- (4.588, 3.419);
	\draw[-latex] (10, 3.75) -- (9, 3.75);
	\node[shape=rectangle, minimum width=0.465cm, minimum height=0.215cm](N1) at (9.5, 4){} node[anchor=center] at (N1.text){$I_2$};
	\node[shape=rectangle, minimum width=0.715cm, minimum height=0.465cm](N2) at (8.625, 2.375){} node[anchor=center] at (N2.text){$V_2$};
	\node[shape=rectangle, xscale=-1, yscale=-1, minimum width=0.465cm, minimum height=0.59cm](N3) at (4.088, 3.669){} node[anchor=center] at (N3.text){$I_1$};
	\node[shape=rectangle, xscale=-1, yscale=-1, minimum width=0.465cm, minimum height=0.465cm](N4) at (4.588, 2.169){} node[anchor=center] at (N4.text){$V_1$};
	\node[shape=rectangle, xscale=-1, yscale=-1, minimum width=0.715cm, minimum height=0.465cm](N5) at (11.625, 2.375){} node[anchor=center] at (N5.text){$V_S$};
\end{tikzpicture}
\caption{Segunda configuración utilizada, generando un circuito abierto en el puerto 1.}
\label{fig:conf_2}


\end{figure}

Luego, para obtener los parámetros restantes, se necesitaba generar un cortocircuito. Sin embargo, por razones de medición, se optó por utilizar una resistencia de valor pequeño, $R_{aux}$, en el puerto donde debería ubicarse el cortocircuito, y calcular la corriente de cortocircuito como la corriente que circula por dicha resistencia.

Se buscó generar un cortocircuito en el puerto 2 usando el circuito que se puede ver en la Figura \ref{fig:conf_3}.

\begin{figure}[H]
\centering
\begin{tikzpicture}[transform shape]
	% Paths, nodes and wires:
	\node[fourport, xscale=1.6, yscale=1.6] at (6.41, 2.251){};
	\node[fourport, xscale=0.3, yscale=0.3] at (5.477, 3.136){};
	\draw (2.25, 3.25) to[american resistor, l_={$R_f$}] (5, 3.25);
	\draw (2.25, 3.25) to[vco] (2.25, 1.25);
	\draw (2.25, 1.25) -- (5, 1.25);
	\draw (7.866, 3.229) -- (9.616, 3.229);
	\draw (7.866, 1.229) -- (9.616, 1.229);
	\draw[-latex] (8.116, 1.479) -- (8.116, 2.979);
	\draw[-latex] (4.75, 1.5) -- (4.75, 3);
	\draw[-latex] (3, 3.812) -- (4, 3.812);
	\draw[-latex] (9.366, 3.479) -- (8.366, 3.479);
	\node[shape=rectangle, minimum width=0.465cm, minimum height=0.215cm](N1) at (8.866, 3.729){} node[anchor=center] at (N1.text){$I_2$};
	\node[shape=rectangle, minimum width=0.715cm, minimum height=0.465cm](N2) at (8.491, 2.354){} node[anchor=center] at (N2.text){$V_2$};
	\node[shape=rectangle, xscale=-1, yscale=-1, minimum width=0.465cm, minimum height=0.59cm](N3) at (3.5, 4.062){} node[anchor=center] at (N3.text){$I_1$};
	\node[shape=rectangle, xscale=-1, yscale=-1, minimum width=0.465cm, minimum height=0.465cm](N4) at (4.5, 2.375){} node[anchor=center] at (N4.text){$V_1$};
	\node[shape=rectangle, minimum width=0.715cm, minimum height=0.465cm](N5) at (1.375, 2.375){} node[anchor=center] at (N5.text){$V_S$};
	\draw (9.616, 3.229) to[american resistor, mirror, l={$R_{aux}$}] (9.616, 1.229);
\end{tikzpicture}

\caption{Tercera configuración utilizada, simulando un cortocircuito en el puerto 2.}
\label{fig:conf_3}
\end{figure}

Por último, se buscó generar un cortocircuito en el puerto 1, como se puede ver en la Figura \ref{fig:conf_4}.

\begin{figure}[H]
\centering
\begin{tikzpicture}[transform shape]
	% Paths, nodes and wires:
	\node[fourport, xscale=1.6, yscale=1.6] at (6.544, 2.272){};
	\node[fourport, xscale=0.3, yscale=0.3] at (5.611, 3.158){};
	\draw (10.75, 3.25) to[american resistor, mirror, l={$R_f$}] (8, 3.25);
	\draw (10.75, 3.25) to[vco, mirror] (10.75, 1.25);
	\draw (8, 1.25) -- (10.75, 1.25);
	\draw (3.338, 3.169) -- (5.088, 3.169);
	\draw (3.338, 1.169) -- (5.088, 1.169);
	\draw[-latex] (8.25, 1.5) -- (8.25, 3);
	\draw[-latex] (4.838, 1.419) -- (4.838, 2.919);
	\draw[-latex] (3.588, 3.419) -- (4.588, 3.419);
	\draw[-latex] (10, 3.75) -- (9, 3.75);
	\node[shape=rectangle, minimum width=0.465cm, minimum height=0.215cm](N1) at (9.5, 4){} node[anchor=center] at (N1.text){$I_2$};
	\node[shape=rectangle, minimum width=0.715cm, minimum height=0.465cm](N2) at (8.625, 2.375){} node[anchor=center] at (N2.text){$V_2$};
	\node[shape=rectangle, xscale=-1, yscale=-1, minimum width=0.465cm, minimum height=0.59cm](N3) at (4.088, 3.669){} node[anchor=center] at (N3.text){$I_1$};
	\node[shape=rectangle, xscale=-1, yscale=-1, minimum width=0.465cm, minimum height=0.465cm](N4) at (4.588, 2.169){} node[anchor=center] at (N4.text){$V_1$};
	\node[shape=rectangle, xscale=-1, yscale=-1, minimum width=0.715cm, minimum height=0.465cm](N5) at (11.625, 2.375){} node[anchor=center] at (N5.text){$V_S$};
	\draw (3.338, 3.169) to[american resistor, l_={$R_{aux}$}] (3.338, 1.169);
\end{tikzpicture}

\caption{Cuarta configuración utilizada, simulando un cortocircuito en el puerto 1.}
\label{fig:conf_4}
\end{figure}

\subsection{Matrices $Z$ experimentales}

Los valores obtenidos para las matrices de impedancias se expresan en $\Omega$.

\[
\mathbf{Z}_{9602} =
\begin{pmatrix}
218.79 \angle -10.96^\circ & 85.98 \angle -67.70^\circ \\
97.24 \angle -72.96^\circ & 157.98 \angle -65.70^\circ
\end{pmatrix}
\ \Omega
\]

\[
\mathbf{Z}_{9610} =
\begin{pmatrix}
326.96 \angle -29.04^\circ & 164.85 \angle -63.24^\circ \\
183,97 \angle -52.04^\circ & 230.08 \angle -40.24^\circ
\end{pmatrix}
\ \Omega
\]

\subsection{Matrices $Y$}

Los valores obtenidos para las matrices de admitancias se expresan en mS.

\subsubsection{Experimentales}

A partir de las mediciones vistas previamente se obtuvieron las matrices:

\[
\mathbf{Y}_{9602} =
\begin{pmatrix}
4.73 \angle 61.60^\circ & 10.72 \angle -1.37^\circ \\
39.46 \angle -13.00^\circ & 6.68 \angle -58.78^\circ
\end{pmatrix}
\ \mathrm{mS}
\]

\[
\mathbf{Y}_{9610} =
\begin{pmatrix}
3.78 \angle 6.32^\circ & 9.00 \angle -1.68^\circ \\
4.08 \angle -27.00^\circ & 4.92 \angle -16.82^\circ
\end{pmatrix}
\ \mathrm{mS}
\]


\subsubsection{Calculadas}

A partir de las matrices $Z$ experimentales, las matrices de admitancias correspondientes se pueden obtener mediante

\begin{equation}
\mathbf{Y}_{Z}=\mathbf{Z}^{-1}.
\end{equation}

Finalmente, la matriz $Y$ obtenida es:

\subsection{Matrices $T$}

Los parámetros de transmisión se presentan en el orden convencional

\begin{equation}
\mathbf{T}=
\begin{bmatrix}
A & B\
C & D
\end{bmatrix}.
\end{equation}

Por lo tanto, el parámetro $B$ posee unidades de impedancia, el parámetro $C$ posee unidades de admitancia, y los parámetros $A$ y $D$ son adimensionales.

\subsubsection{Experimentales}

A partir de las mediciones vistas previamente se obtuvieron las matrices:

\[
\mathbf{T}_{9602} =
\begin{pmatrix}
2.25 \angle 62.00^\circ & -25.34 \angle 13.00^\circ \ \Omega \\
10.28 \angle 72.96^\circ \ \mathrm{mS} & -0.120 \angle 74.60^\circ
\end{pmatrix}
\]

\[
\mathbf{T}_{9610} =
\begin{pmatrix}
16.62 \angle 23.00^\circ & -245.00 \angle 27.00^\circ \ \Omega \\
50.83 \angle 52.04^\circ \ \mathrm{mS} & -0.926 \angle 33.32^\circ
\end{pmatrix}
\]


\subsubsection{Calculadas}

Del mismo modo, utilizando la convención adoptada para los parámetros de transmisión, se utilizaron las siguientes ecuaciones para obtener la matriz $T$ a partir de la matriz $Z$:

\begin{equation}
A=\frac{Z_{11}}{Z_{21}},
\qquad
B=\frac{\det(\mathbf{Z})}{Z_{21}},
\qquad
C=\frac{1}{Z_{21}},
\qquad
D=\frac{Z_{22}}{Z_{21}}.
\end{equation}

Obteniendo así las matrices:
